% Paper 3, §13 — Frontier-2026 snapshot: reasoning protocols
% Anchors: kb/index/timeline.md 2024–2026 entries
%          kb/notes/scaling/inference-time-compute-scaling.md §3.3
%          kb/excerpts/{deepseek-r1,s1-2025}.md plus model-card papers
% Cross-paper: Paper-1 §13 is the architecture-side analog; this is the
%              reasoning-protocol side.

\section{Frontier-2026 snapshot: reasoning protocols}
\label{sec:frontier-2026}
\label{sec:frontier-snapshot}

The previous twelve sections decompose reasoning into three legs ---
inference-time compute (\cref{sec:inference-compute-scaling,sec:self-consistency,sec:budget-forcing}), verifier signals
(\cref{sec:process-reward-models,sec:tree-search,sec:prm-survival}), and
RL on verifiable rewards
(\cref{sec:rlvr-grpo-dapo,sec:r1-family}) --- and ask, leg by leg, what
the literature settled and what remains contested. This section closes
that decomposition by collapsing it into one cross-cutting view: which
\emph{currently-deployed} reasoning model ships which legs of the
triangle, and at what budget. The intended use is the same as Paper~1's
architecture-side frontier snapshot --- a reference figure that lets a
reader place any new model release on the design space the rest of the
paper formalises. Because the underlying training and serving stacks
move on the order of weeks, this snapshot is dated 2026-Q2 and explicit
about which entries are tier-A primary sources, which are vendor system
cards, and which are still discovery-tier signals pending paper-grade
documentation.

The argument structure is: a single 2-page comparison table
(\cref{tab:frontier-2026}) that codes each model along the three legs;
three short subsections grouping the lineup into R1-style, o-series-style,
and hybrid-think families; an agreement-and-divergence summary; and a
forward link into the contradictions catalogued by \cref{sec:contradictions}.
The single load-bearing claim of the section is structural rather than
empirical: \emph{every} frontier-2026 reasoning entry ships RLVR plus
long-CoT, \emph{none} of them ship a process reward model in the headline
RL loop, and the productionised exposure of inference-time compute has
converged on a small number of API-shaped ``thinking budget'' surfaces.

\subsection{The 2-page comparison table}
\label{sec:frontier-table}

\Cref{tab:frontier-2026} summarises six reasoning entries that
collectively span the open-weight, open-API, and closed-frontier ends of
the 2026 lineup. The columns trace the three legs of the triangle plus
two productionisation columns. The legend below the table specifies
exactly what counts as ``ships'' for each leg.

\begin{table}[h]
\centering
\footnotesize
\renewcommand{\arraystretch}{1.25}
\begin{tabularx}{\textwidth}{@{}lXXXXX@{}}
\toprule
\textbf{Model} & \textbf{Long-CoT} & \textbf{RLVR} & \textbf{Verifier in RL loop} & \textbf{Inference search} & \textbf{Budget exposure} \\
\midrule
DeepSeek-R1 \citep{deepseek-r1}
  & yes; ${\sim}10^4$ tok traces \citep[\S 2.3]{deepseek-r1}
  & yes; GRPO \citep[\S 2.1]{deepseek-r1}, terminal binary reward
  & none; explicit no-PRM commitment \citep[\S 2.2]{deepseek-r1}
  & none in headline; self-consistency optional \citep[\S 2.3]{deepseek-r1}
  & none (open-weights, fixed schema) \\
Qwen3-thinking
  & yes; toggleable via system prompt
  & yes \deepencite{qwen3-thinking — system card pending in references.bib; cf.\ \citet{qwen3} for the lineage and base}
  & verifier-only in headline reports \deepencite{qwen3-thinking}
  & none in headline
  & think / non-think mode toggle \deepencite{qwen3-thinking} \\
QwQ
  & yes; reasoning-preview model
  & yes \deepencite{qwq — Qwen's standalone reasoning preview, predates Qwen3-thinking; system card pending}
  & verifier-only \deepencite{qwq}
  & none in headline
  & none (always-on thinking) \\
OpenAI o1 / o3
  & yes; thinking-first generation
  & yes (RL-class on reasoning trajectories per the system card; whether the reward is verifier-derived in the RLVR sense is not disclosed): ``the more reinforcement learning ($\ldots$) the more time it spends thinking'' \deepencite{openai-o1 — system card; cf.\ openai-o3 — system card; claimed in the system card, not externally reproduced}
  & not disclosed; headline reports do not name a PRM \deepencite{openai-o1}
  & not disclosed
  & \texttt{reasoning\_effort} $\in\{\text{low, medium, high}\}$ \deepencite{openai-o-series — API reference}\\
Gemini~2.5-thinking \citep{gemini2-5}
  & yes; reasoning-mode
  & yes \deepencite{gemini2-5 — Google DeepMind technical report}
  & not disclosed; verifier-style reward emphasised \deepencite{gemini2-5}
  & not disclosed
  & \texttt{thinkingBudget} (token cap) \deepencite{gemini2-5 — API reference} \\
Claude 3.7 / 4 extended thinking
  & yes; ``extended thinking'' mode
  & yes \deepencite{claude-3-7 — Anthropic system card; analogous treatment for Claude 4 family}
  & not disclosed; constitutional + verifier-style mix \deepencite{claude-3-7}
  & not disclosed
  & extended-thinking toggle + budget knob \deepencite{claude-3-7} \\
\bottomrule
\end{tabularx}
\caption{Frontier-2026 reasoning lineup along the three legs of the
inference-scaling triangle. ``Long-CoT'': the deployed protocol emits a
\texttt{<think>...</think>}-style trace structurally distinct from the
answer phase, with the model trained to use the channel
\citep[\S 2]{deepseek-r1}.\protect\footnotemark{} ``RLVR'': the
post-training pipeline includes an RL phase whose reward depends on a
mechanically-checkable correctness signal (math answer, unit-test pass,
regex), as formalised by \cref{eq:verifier-reward} of
\cref{sec:rlvr-verifier-reward}. ``Verifier in RL loop'': whether a
\emph{neural} reward model (PRM or ORM) appears alongside the verifiable
reward during RL, as opposed to terminal binary verifiers only.
``Inference search'': whether the deployed inference path explicitly
expands a beam, MCTS tree, or PRM-guided rollout (vs.\ pure long-CoT or
$N$-sample self-consistency on top of it). ``Budget exposure'': the API
or system-prompt surface through which a caller can trade thinking-token
spend for accuracy at request time, the productionised form of Snell's
\cref{eq:snell-compute-optimal} \cref{sec:snell-compute-optimal}.
Categorisation rationale per row in
\cref{sec:frontier-r1-family,sec:frontier-o-series,sec:frontier-hybrid-budget}.
Open-weight rows cite tier-A primary sources; closed-weight rows cite
tier-B vendor system cards via \deepencite{} where \texttt{papers.json}
does not yet carry a stable bibkey.}
\label{tab:frontier-2026}
\end{table}
\footnotetext{KB excerpt: \texttt{kb/excerpts/deepseek-r1.md\#sec-2}.}

Three reading conventions follow from the table. First, the
``Verifier in RL loop'' column is uniformly ``none'' or ``not
disclosed'': no frontier-2026 entry that has published a system card
identifies a PRM as a reward source in the headline RL loop. The R1
authors' explicit no-PRM commitment
\citep[\S 2.2]{deepseek-r1}\footnote{KB excerpt:
\texttt{kb/excerpts/deepseek-r1.md\#sec-2-2-reward}.} is the strongest
attested form; the closed-weight entries decline to specify, which on
this question is operationally indistinguishable from absence at the
level of evidence available.

Second, the ``Inference search'' column is also uniformly ``none in
headline'' or ``not disclosed.'' The MCTS-and-PRM-guided lineage of
\cref{sec:tree-search} is a research lineage; deployed reasoning models
in 2026 ship long-CoT plus optional self-consistency on top. Whether
that is because search does not pay at frontier scale, because the
extra orchestration is a serving-cost burden, or because the PRMs that
would guide it are not strong enough on reasoning-trained policies
(\cref{sec:prm-survival}) is not separable from the published 2026
evidence base.

Third, the ``Budget exposure'' column is the column where the field
visibly converges on an API shape: a discrete-level reasoning-effort
selector \deepencite{openai-o-series — API reference}, a token-cap
\texttt{thinkingBudget} \deepencite{gemini2-5 — API reference}, an
extended-thinking toggle \deepencite{claude-3-7}, or a system-prompt
think / non-think mode \deepencite{qwen3-thinking}. All four are the
productionised form of $\theta$ in \cref{eq:snell-compute-optimal} ---
the strategy choice the section-4 anchor formalises as the
prompt-conditional argmax over compute spend.

\subsection{R1-style: long-CoT + RLVR + verifier-only}
\label{sec:frontier-r1-family}

The R1 family is the cleanest published form of the triangle: open
weights, an open paper with explicit hyperparameters, a verifier-only
RL loop \citep[\S 2.2]{deepseek-r1}, and an RL phase that drives
average response length from $\sim\!200$ to $\sim\!10^4$ tokens with
no length term in the reward
\citep[\S 2.3]{deepseek-r1}.\footnote{KB excerpt:
\texttt{kb/excerpts/deepseek-r1.md\#sec-2-3-aha}.} \Cref{sec:r1-family}
walked the four-stage protocol; the family it seeded by 2026-Q2
includes Qwen's QwQ reasoning-preview \deepencite{qwq} and Qwen3-thinking
\deepencite{qwen3-thinking; cf.\ \citet{qwen3}}, plus the open-replication
track Tülu~3 \citep{tulu3} and the Hugging Face OpenR1 community effort
\deepencite{openr1 — community reproduction; pending bibkey}.

The structural commitment shared across this row of
\cref{tab:frontier-2026} is doctrinal more than algorithmic: a binary
verifier reward in \cref{eq:verifier-reward} is treated as the gold
reward in its supported domains, and a learned reward model is treated
as a reward-hackable proxy that the field has now opted out of for the
reasoning subset of post-training. \citet[\S 2.2]{deepseek-r1} state
the position; the open-replication track inherits it; the closed-frontier
rows of the table neither contradict it on the published evidence nor
positively confirm it.

\subsection{o-series: RL-trained thinking with proprietary mix}
\label{sec:frontier-o-series}

OpenAI's o1 (Sept 2024) and o3 successor lineage \deepencite{openai-o1
— system card; openai-o3 — system card and capability writeups} are
treated by the published material as the productionised form of the
same ``RL on reasoning trajectories'' family that R1 reproduces in the
open. The system-card narrative is uniformly that ``the more
reinforcement learning ($\ldots$) and the more time it spends thinking,
the better it performs on reasoning tasks''
\deepencite{openai-o1 — system card}; no PRM is named in the headline
reward source; no inference-time search algorithm is named in the
deployed serving path. The o-series exposes a discrete
\texttt{reasoning\_effort} parameter at the API level
\deepencite{openai-o-series — API reference} that maps to a
thinking-token-spend budget, the same operative knob as Gemini~2.5's
\texttt{thinkingBudget} but discretised rather than token-quantised.

The reading available from the published material: the o-series and the
R1 family are operationally close enough on the published axes
(long-CoT yes, RLVR yes, terminal verifier in the publicly-disclosed
reward, no headline PRM, no headline inference search) that the
remaining differences are training-data composition, base-model scale,
and proprietary post-training tooling. None of those three are visible
in the 2026 system-card disclosure level. The cross-paper hook to
Paper~1's architecture-side snapshot is symmetric: closed-frontier
architecture rows are also dominated by ``not disclosed'' on the
load-bearing axes, and the convergence-vs-divergence reading in this
paper depends on the same disclosure ceiling.

\subsection{Hybrid think / non-think with productionised budget knobs}
\label{sec:frontier-hybrid-budget}

The third 2026 productionisation pattern --- and the operationally
distinct one --- is the hybrid model that ships \emph{both} a
non-thinking decode path and a thinking-decode path, switched at
request time by a system prompt or API parameter. Three exemplars:

\begin{itemize}
\item \textbf{Gemini~2.5} \citep{gemini2-5} exposes a
\texttt{thinkingBudget} parameter that maps to a hard cap on
thinking-token spend per request, with cheap mode and expensive mode at
the two ends of the budget axis \deepencite{gemini2-5 — API reference}.
This is the most-explicit token-cap formulation of
\cref{eq:snell-compute-optimal} in deployment.
\item \textbf{Anthropic Claude 3.7 (and the Claude~4 family)} exposes
``extended thinking'' as a request-time toggle alongside a budget knob
\deepencite{claude-3-7 — Anthropic system card; analogous treatment in
the Claude~4 family}. The non-extended-thinking decode path remains
production-standard for latency-sensitive consumer endpoints; the
extended-thinking path is invoked when the caller has a hard problem.
\item \textbf{Qwen3-thinking} \deepencite{qwen3-thinking — system card}
is the open-weight version of the same pattern: a single model that
toggles between thinking and non-thinking modes via a system-prompt
flag, with mode-tagged post-training that shapes both branches.
\end{itemize}

The shared structural commitment across the hybrid row of
\cref{tab:frontier-2026} is that the productionisation surface
acknowledges \cref{eq:snell-compute-optimal} explicitly: the caller
\emph{is} the strategy chooser $\theta^\ast$, and the API exposes the
budget axis $N$ as a first-class request-time parameter. This is the
operational form of the difficulty-binning operationalisation of
\citet[\S 3.2]{snell2024}\footnote{KB excerpt:
\texttt{kb/excerpts/snell2024.md\#sec-3-2-difficulty}.} ---
hard-problem callers select expensive mode; easy-problem callers
select cheap mode --- pushed to the API edge.

\begin{intuition}
The hybrid budget knob is Snell's compute-optimal selection problem
shipped to the caller. \Cref{eq:snell-compute-optimal} defines
$\theta^\ast_{q, y^\ast(q)}(N)$ as a function of the prompt $q$;
the deployed API treats $\theta$ as a parameter the \emph{caller}
sets per request, on the assumption that the caller has more
information about the problem's difficulty than a difficulty-classifier
that the model would otherwise have to learn. Returning to math: the
expectation
$\E_{y \sim \mathrm{Target}(\theta, N, q)}[\mathbb{1}_{y = y^\ast(q)}]$
that \cref{sec:snell-compute-optimal} maximises is now maximised by a
two-player decomposition --- the model fixes the policy
$\pi_\theta(\cdot \mid q;\, \text{thinking budget } L)$, and the
caller fixes $L$ at request time. The compute-optimal $\theta^\ast(N)$
of Snell's formal statement is then the joint optimum the two players
together approximate, with the API surface as the coordination
mechanism.
\end{intuition}

\subsection{Where the lineup agrees}
\label{sec:frontier-agree}

Four cross-row agreements are stable across the 2026-Q2 published
material:

\begin{enumerate}
\item \textbf{RLVR is universal at frontier.} Every reasoning row of
\cref{tab:frontier-2026} reports an RL phase whose reward signal
includes a verifiable correctness component on math, code, or both.
The R1 row is the most explicit \citep[\S 2.2]{deepseek-r1}; the
o-series and Gemini~2.5 rows attribute headline gains to ``more
reinforcement learning'' \deepencite{openai-o1; gemini2-5} on
reasoning trajectories. The DAPO-class RL improvements
of \cref{sec:rlvr-grpo-dapo} \citep{dapo2025}\footnote{KB excerpt:
\texttt{kb/excerpts/dapo2025.md\#sec-2}.} are part of the same family
deployed in some open-weight reproductions \citep{tulu3}.

\item \textbf{Long-CoT is universal.} Every row deploys a
\texttt{<think>...</think>}-style structurally-separate reasoning
phase ahead of the answer. The pattern is no longer paper-grade
research; it is the deployment standard. \Cref{sec:r1-family} and
\cref{sec:budget-forcing} cover the training-time and inference-time
mechanisms, respectively.

\item \textbf{PRMs are absent from headline RL loops.} The
``Verifier in RL loop'' column is empty in every row that has
published a system card, and explicitly empty in the R1 row
\citep[\S 2.2]{deepseek-r1}. The post-R1 question
\cref{sec:prm-survival} catalogues --- whether PRMs survive as a
\emph{training} signal beyond verifier-only RLVR --- is not answered
by current frontier deployments; it is answered by their non-adoption.

\item \textbf{Thinking-budget exposure is the productionised form of
the inference-compute axis.} Every row that exposes any
inference-compute knob does so as a budget cap or discrete-effort
parameter --- the operational form of $\theta$ in
\cref{eq:snell-compute-optimal}. The convergence on this surface is
strong enough that a 2026 caller can transfer expectations about
what ``high reasoning effort'' means across vendors, even though the
underlying policies and base models are distinct.
\end{enumerate}

\subsection{Where the lineup diverges}
\label{sec:frontier-diverge}

Three load-bearing axes of divergence remain, all of them either not
disclosed by closed-frontier vendors or actively contested in the
open-replication literature.

\paragraph{Training-data composition.} R1's pretraining mix is
DeepSeek-V3's ($\sim\!14.8\text{T}$ tokens, mix not fully disclosed
\citep[\S 2]{deepseek-r1}); the o-series and Claude families do not
publish their pretraining mixes; Gemini~2.5 publishes only summary
characterisations \deepencite{gemini2-5}. The AIME-2024 contamination
contradiction of \cref{sec:r1-contamination} interacts non-trivially
with this axis: rows whose pretraining mix is undisclosed cannot be
contamination-audited from outside, and the headline reasoning numbers
on AIME-class benchmarks are conditional on the audit not having been
done.

\paragraph{Reasoning-architecture details.} Whether the model is a
single dense Transformer trained with mode-tagged data
\deepencite{qwen3-thinking}, a model with a separate thinking-decode
path \deepencite{claude-3-7 extended thinking}, or a model with no
public architectural documentation at all (the o-series) is not
visible to a caller, and is operationally distinct from the
training-protocol differences of \cref{sec:r1-fourstage}. Paper~1's
reasoning-architecture axis is the architecture-side counterpart;
this paper's row taxonomy in \cref{tab:frontier-2026} is the
training-protocol-side counterpart.

\paragraph{Search vs.\ no-search at inference.} The ``Inference
search'' column is uniformly ``not disclosed'' or ``none.'' Whether
some rows in fact run undisclosed PRM-guided beam or MCTS rollouts in
serving is not separable from the available evidence. The research
literature \cref{sec:tree-search} continues to push search-and-PRM
methods on benchmarks; the deployed lineup either does not adopt them
or does not disclose adoption. A correct 2026-Q2 reading is that
search at inference is currently a research-and-eval-stack lineage,
not a deployment-stack lineage, with the deployment-stack reasons
either unrelated to algorithmic value or undisclosed.

\begin{contradiction}
\textbf{The disclosure ceiling and what it means for the snapshot.}
\Cref{tab:frontier-2026} reports the 2026-Q2 published material at
face value: \emph{published} system cards say ``no PRM in headline
reward,'' \emph{published} reports describe long-CoT and RLVR, and
\emph{published} APIs expose budget knobs. None of those statements
rule out the contrary on closed-weight rows. The R1 row is the only
one where the ``Verifier in RL loop = none'' cell is positively
attested rather than not-disclosed
\citep[\S 2.2]{deepseek-r1}.\footnote{KB excerpt:
\texttt{kb/excerpts/deepseek-r1.md\#sec-2-2-reward}.} Two readings of
the column-uniformity that follows are both compatible with the
evidence: \emph{(i)} the field has converged on verifier-only RL plus
long-CoT plus thinking-budget knobs as the dominant pattern, or
\emph{(ii)} the closed-frontier vendors are running undisclosed PRM
or search machinery and the disclosure ceiling produces the apparent
convergence as a measurement artefact. The 2026-Q2 published evidence
does not discriminate; \cref{sec:contradictions} catalogues this as a
load-bearing open contradiction whose resolution requires either
disclosure changes or third-party reverse-engineering at frontier
scale.
\end{contradiction}

\subsection{Cross-paper hook and forward link}
\label{sec:frontier-forward}

Paper~1's frontier-2026 architecture snapshot covers the
architecture-side analog of \cref{tab:frontier-2026}: convergence on
KV-sharing variants (GQA / MLA), RMSNorm with Pre-LN, RoPE with
position-extension methods, SwiGLU FFNs, MoE at scale. The reasoning
side of the same lineup --- this section's table --- exhibits the same
structural pattern: a small number of dominant-deployment choices
(verifier-only RLVR, long-CoT trace, thinking-budget API), a small
number of contested axes (training-data composition, search-at-inference,
reasoning-architecture details), and a long tail of vendor-specific
proprietary tooling that the 2026 disclosure regime does not surface.
The two snapshots together are the synthesis the rest of the series'
section-by-section treatments build toward.

The forward link is to \cref{sec:contradictions}, which catalogues the
six load-bearing open contradictions of the paper:
PRM utility post-R1 (\cref{sec:prm-survival}), CoT faithfulness
scaling (\cref{sec:cot-faithfulness}), RLVR generalisation versus
memorisation
(\cref{sec:rlvr-grpo-dapo,sec:r1-family}), search-vs-RL under fixed
compute (\cref{sec:search-vs-rl}), AIME-2024 contamination
(\cref{sec:r1-contamination}), and MCTS-vs-beam at matched budget
(\cref{sec:tree-search}). The frontier-2026 snapshot is not a
resolution of those contradictions --- it is a coordinate system in
which they are localisable. A reader who knows where a 2026 model sits
in \cref{tab:frontier-2026} can read off which of the six
contradictions bears most directly on its expected behaviour, and which
of the disclosed and undisclosed axes most affect what a third party
can verify about it. The headline of the snapshot is structural and
sober: the triangle is real, the three legs co-design, the
productionisation surface has converged, and the open questions about
which combination is optimal at which (compute, latency, domain)
operating point will not be closed-form for at least another model
generation.
